\documentclass[11pt]{article}
\usepackage[T1]{fontenc}
\usepackage[utf8]{inputenc}
\usepackage{lmodern}
\usepackage[margin=1in]{geometry}
\usepackage{microtype}
\usepackage{amsmath,amssymb,amsthm,mathtools}
\usepackage{enumitem}
\usepackage{xcolor}
\usepackage{hyperref}
\usepackage{booktabs}
\usepackage{array}
\usepackage{tabularx}
\usepackage{mdframed}
\hypersetup{colorlinks=true,linkcolor=blue!45!black,citecolor=blue!45!black,urlcolor=blue!45!black}

\definecolor{accent}{RGB}{34,71,120}
\definecolor{lightaccent}{RGB}{236,242,249}

\setlength{\parindent}{0pt}
\setlength{\parskip}{0.55em}
\setlist[itemize]{topsep=0.25em,itemsep=0.2em,leftmargin=1.5em}
\setlist[enumerate]{topsep=0.25em,itemsep=0.2em,leftmargin=1.8em}

\newtheorem{proposition}{Proposition}[section]
\newtheorem{remark}[proposition]{Remark}
\newtheorem{definition}[proposition]{Definition}

\newcommand{\ALCIRCC}{\mathsf{ALCI}_{\mathsf{RCC5}}}
\newcommand{\Eq}{\mathsf{EQ}}
\newcommand{\Dr}{\mathsf{DR}}
\newcommand{\Po}{\mathsf{PO}}
\newcommand{\Pp}{\mathsf{PP}}
\newcommand{\Ppi}{\mathsf{PPI}}
\newcommand{\FW}{\mathsf{FW}}
\newcommand{\md}{\operatorname{md}}
\newcommand{\cl}{\operatorname{cl}}

\title{\textbf{What Is Still Missing for Decidability of $\ALCIRCC$?}\\[0.35em]
\large After the failure of fixed-width extraction}
\author{Prepared by GPT-5.4 Thinking}
\date{31 March 2026}

\begin{document}
\maketitle

\begin{mdframed}[backgroundcolor=lightaccent,linecolor=accent]
\textbf{Purpose.}
This short note integrates three strands of the current discussion around $\ALCIRCC$:
Wessel's original open-problem report \cite{Wessel2003}, the later quasimodel draft \cite{Claude2026Quasi}, and the new counterexample showing that the fixed-width extraction conjecture $\FW(C,N)$ fails \cite{Claude2026FW}. It also records what still survives from the contextual-signature counting approach developed in the accompanying 2026 notes \cite{GPT2026Bound,GPT2026Tableau}.
\end{mdframed}

\begin{abstract}
The decidability of $\ALCIRCC$ remains open. Two proposed 2026 routes are now blocked for different reasons. First, the quasimodel paper does not yet yield a correct completeness proof: the abstraction by unary types and pair-types is too coarse, and the patchwork step does not justify the needed global amalgamation \cite{Claude2026Quasi}. Second, the contextual tableau route is no longer merely incomplete: the attached counterexample proves that its key finite-width extraction conjecture $\FW(C,N)$ is false for every $N$ on the satisfiable concept
\[
C_\infty=(\exists \Pp.\top)\sqcap(\forall \Pp.\exists \Pp.\top),
\]
which has only infinite models \cite{Claude2026FW,Wessel2003}.

What survives is a weaker local result: around a fixed finite designated core, one can still bound the number of bounded-depth contextual signatures by a computable function of $|C|$ \cite{GPT2026Bound}. What is missing now is a theorem that converts this local finiteness into a genuine decision procedure. In concrete terms, one needs a \emph{regular omega-model theorem}: every satisfiable concept should admit a representation by finitely many local interface signatures together with finitely many abstract $\Pp/\Ppi$ thread states and an acceptance condition handling infinite proper-part chains. Without such an omega-level regularity theorem, the current finite abstractions are either too weak (unary quasimodels) or too rigid (fixed-width exact recentering).
\end{abstract}

\section{Background}

Wessel's report introduced the $\mathsf{ALCIRCC}$ family and explicitly left the decidability of $\ALCIRCC$ and $\mathsf{ALCI}_{\mathsf{RCC8}}$ open. The report also emphasizes the two structural obstacles that distinguish these logics from ordinary description logics: models are complete graphs rather than trees, and the concept
\[
(\exists \Pp.\top)\sqcap(\forall \Pp.\exists \Pp.\top)
\]
shows that there is no finite model property \cite[pp.~16--20, 48--50]{Wessel2003}. Any successful decision procedure therefore has to cope with infinite $\Pp$-chains directly.

The later quasimodel draft claims positive decidability results for $\ALCIRCC$ and $\mathsf{ALCI}_{\mathsf{RCC8}}$ via a type-based abstraction and patchwork arguments \cite{Claude2026Quasi}. However, that route still has serious proof problems: the abstraction identifies elements by unary closure types and pair-types, which loses witness identity, and the argument from local compatibility to a single global assignment is not justified by the RCC patchwork property in the form used there.

In response, the contextual tableau idea replaced plain unary types by finite contextual local states. Its soundness direction is plausible: an open tableau graph can be unfolded into a genuine model. Completeness, however, was reduced to a finite-width extraction conjecture $\FW(C,N)$ stating that every satisfiable concept has some bounded-width closed family of such local states \cite{GPT2026Tableau}.

\section{What the new counterexample settles}

The attached note proves that $\FW(C,N)$ is false for every $N$ on the concept
\[
C_\infty=(\exists \Pp.\top)\sqcap(\forall \Pp.\exists \Pp.\top).
\]
The mechanism is simple and convincing \cite{Claude2026FW}:
\begin{itemize}
\item In any local state, the relation $\Pp$ on the remembered items is a finite strict partial order.
\item Every finite strict partial order has a maximal element.
\item But the recentering axioms force every designated $\Pp$-witness to become the center of a successor state whose own $\exists\Pp.\top$ demand must again be witnessed.
\item When those successor witnesses are transported back into the original root state, one obtains an arbitrarily long strictly ascending $\Pp$-chain inside one finite carrier, contradiction.
\end{itemize}

This is not a refutation of satisfiability of $C_\infty$; indeed $C_\infty$ has obvious infinite chain models. It is a refutation of the claim that such models can always be compressed into a \emph{finite closed family of exact local states with total recentering back into a finite parent carrier}. That is the crucial point.

\begin{proposition}
After \cite{Claude2026FW}, the fixed-width contextual tableau strategy is no longer an incomplete proof attempt. Its central completeness conjecture is false.
\end{proposition}

\section{How this interacts with the earlier 2026 notes}

The main integration point is a distinction between two different finiteness claims.

\subsection*{(A) Strong finiteness: exact local-state closure}
This is the claim attacked by $\FW(C,N)$. It says, in effect, that every satisfiable concept can be represented by finitely many exact local states whose successor states recenter back into one another via literal item-by-item maps. The new counterexample shows that this is too strong.

\subsection*{(B) Weak finiteness: bounded local descriptors}
By contrast, the finite-width bound from the signature-counting note only asserts that, around a fixed finite designated core, there are only finitely many \emph{bounded-depth contextual signatures} relevant to formulas of modal depth at most $\md(C)$ \cite{GPT2026Bound}. This is a local counting statement, not a claim that all those signatures can be assembled into a finite exact closed family obeying the recentering axioms.

So the new counterexample does \emph{not} automatically destroy the signature-counting result. What it destroys is the stronger step that tried to pass from local finite descriptors to a closed finite-state tableau object with exact recentering. The correct lesson is:
\begin{quote}
Local descriptor finiteness may still be true and useful, but it is not enough on its own. Infinite $\Pp$-progress cannot be encoded by exact finite recentering.
\end{quote}

\section{Current status of the three routes}

\begin{center}
\renewcommand{\arraystretch}{1.25}
\begin{tabularx}{\textwidth}{@{}>{\raggedright\arraybackslash}p{0.26\textwidth}>{\raggedright\arraybackslash}p{0.26\textwidth}>{\raggedright\arraybackslash}X@{}}
\toprule
\textbf{Route} & \textbf{What works} & \textbf{Why it does not yet decide $\ALCIRCC$} \\
\midrule
Wessel 2003 \cite{Wessel2003} & Correct problem analysis; lower bounds; explicit identification of infinite $\Pp$-chains and lack of tree/finite model properties. & Leaves decidability open. \\
\addlinespace
Quasimodel draft \cite{Claude2026Quasi} & The type-level setup is a natural first abstraction. & The proof still loses witness identity and misuses patchwork at the global amalgamation step. \\
\addlinespace
Contextual tableau draft \cite{GPT2026Tableau} & Soundness direction is plausible: an accepted abstract object can be unfolded into a model. & Completeness rested on $\FW(C,N)$, now refuted by \cite{Claude2026FW}. \\
\addlinespace
Contextual-signature bound \cite{GPT2026Bound} & Provides a computable local descriptor bound around a fixed finite core. & Does not explain how to represent infinite $\Pp/\Ppi$ progress by a finite accepted object. \\
\bottomrule
\end{tabularx}
\end{center}

\section{What is now missing}

The missing ingredient is not a small patch. It is a new \emph{regularity theorem} for infinite proper-part progress.

\begin{definition}[informal target theorem]
A \emph{regular omega-model theorem} for $\ALCIRCC$ would say that every satisfiable concept $C$ has a model representable by:
\begin{enumerate}
\item a finite alphabet of local interface signatures derived from $\cl(C)$,
\item finitely many abstract control states for $\Pp/\Ppi$ threads, and
\item an acceptance condition guaranteeing that infinite proper-part chains are realized by fresh progress rather than by illicit reuse inside a finite strict partial order.
\end{enumerate}
\end{definition}

In practice, this means four separate tasks remain.

\begin{enumerate}
\item \textbf{Replace exact recentering by interface transfer.}
A successor state should preserve only the finite interface to the old context, not embed its entire item set back into the parent.

\item \textbf{Add an explicit omega-mechanism.}
A calculus for $C_\infty$ must accept an infinite $\Pp$-thread. Wessel already suggested that something like an ``infinity checker'' would be needed here \cite[p.~49]{Wessel2003}. The new counterexample confirms that ordinary blocking is not enough.

\item \textbf{Prove regular extraction.}
One needs a theorem that every satisfiable concept admits a model whose infinite $\Pp$-behavior is regular at the level of the finite interface signatures.

\item \textbf{Prove realization.}
Given such a finite omega-regular abstract object, one must still show that it expands to a genuine RCC5 model. This is the point where the patchwork property should enter, but only after the abstraction has been designed correctly.
\end{enumerate}

\section{A plausible next route}

The most promising route now seems to be the following hybrid strategy.

\begin{enumerate}
\item Keep the contextual-signature bound as the \emph{finite alphabet}. It still gives a computable repertoire of local interface types \cite{GPT2026Bound}.
\item Abandon exact fixed-width recentering. The counterexample shows that literal finite carrier reuse is the wrong invariant \cite{Claude2026FW}.
\item Replace it with a Buchi- or parity-style acceptance condition for $\Pp/\Ppi$ threads. Infinite proper-part progress should be handled semantically by accepted runs, not by embedding an infinite chain into one finite local state.
\item Then prove a regular-model theorem and a realization theorem.
\end{enumerate}

This is also the route most consistent with Wessel's original diagnosis. The core issue was never merely ``how many local configurations are there?''; it was always ``how can a finite procedure soundly represent infinite $\Pp$-progress in complete-graph models?'' \cite{Wessel2003}. The new counterexample clarifies that the answer cannot be ``by exact finite recentering.''

\section{Conclusion}

The state of play is now clearer than before.

\begin{itemize}
\item The quasimodel proof is still not a finished decidability proof.
\item The fixed-width contextual tableau completeness conjecture is false.
\item A weaker finite local descriptor bound may still be correct and useful.
\item The missing step is now sharply identifiable: a regular omega-model theorem handling infinite $\Pp/\Ppi$ chains.
\end{itemize}

So the decidability of $\ALCIRCC$ remains open, but the open problem is better focused than before. The right target is no longer a finite closed family of exact local states. It is a finite \emph{omega-regular} abstraction of complete-graph models with proper-part threads.

\begin{thebibliography}{9}

\bibitem{Claude2026FW}
Claude (Anthropic).
\newblock \emph{On the Failure of Finite-Width Extraction for $\mathsf{ALCI\_RCC5}$: A Counterexample to the $\mathsf{FW}(C,N)$ Conjecture and Its Implications for the Contextual Tableau Approach}.
\newblock Discussion draft, 2026.

\bibitem{Claude2026Quasi}
Claude (Anthropic).
\newblock \emph{On the Decidability of $\mathsf{ALCIRCC5}$ and $\mathsf{ALCIRCC8}$: Quasimodels Meet the Patchwork Property}.
\newblock Discussion draft, 2026.

\bibitem{GPT2026Bound}
GPT-5.4 Pro.
\newblock \emph{A Finite-Width Bound for Contextual Signatures in $\mathsf{ALCI}_{\mathsf{RCC5}}$}.
\newblock Discussion draft, 2026.

\bibitem{GPT2026Tableau}
GPT-5.4 Pro.
\newblock \emph{A Contextual Tableau Calculus for $\mathsf{ALCI}_{\mathsf{RCC5}}$}.
\newblock Discussion draft, 2026.

\bibitem{Wessel2003}
Michael Wessel.
\newblock \emph{Qualitative Spatial Reasoning with the ALCIRCC Family - First Results and Unanswered Questions}.
\newblock Technical Report FBI-HH-M-324/03, University of Hamburg, 2002/2003.

\end{thebibliography}

\end{document}
